\section{Generic Cubic Two-Matrix Model}
To illustrate our formalism and demonstrate the strength of the loop equation approach, let us consider the simple example of a two-matrix model with a generic cubic potential:
%
\begin{align}
    \label{2MMGeneric}
    V(M_1,M_2) = & \frac{1}{2(\alpha \beta - \gamma^2)}(\alpha M_1^2 + \beta M_2^2 - 2 \gamma M_1 M_2) \\ & \qquad \quad \nonumber + \frac{1}{3} (a M_1^3 + 3b M_1^2 M_2 + 3d M_1 M_2^2 + e M_2^3).
\end{align}
%
This furnishes an example of a multi-matrix model which is not solvable through any known saddle-point or orthogonal polynomial approach. In particular, the loop equation method of Chapter 2 is not sufficient to compute the planar resolvent. However, as first demonstrated in \cite{Carroll:1995nj}, one may compute the resolvent of this model using our more granular loop equation strategy.

For \eqref{2MMGeneric} the following are the minimal set of reparameterisations:
%
\begin{align}
    \label{2MMGenericReps}
    \delta M_{i} & = \frac{1}{z-M_1}, \\
    \delta M_{i} & = M_2 \frac{1}{z-M_1} + \text{h.c.},\\
    \delta M_{i} & = M_2^2 \frac{1}{z-M_1}+\text{h.c.},\\
    \delta M_{i} & = M_2 \frac{1}{z-M_1} M_2,
\end{align}
%
for $i\in\{1,2\}$. Calculating the corresponding loop equations yields the following relations, where the action of the $\Delta$ operator is given by \eqref{delta} for $k=1$, i.e. $\Delta^1=\Delta$, and we suppress the $z$-dependence of the resolvent functions\footnote{Note that $W_2$ is distinct from the resolvent $W_{(2)}$ \eqref{eq2}, which denotes the resolvent of the sum of two matrices.}:
%
\begin{align}
    W^2 =&  -d W_{(2,2)}+\frac{\gamma  W_{2}}{\gamma ^2-\alpha  \beta }-a\Delta^2 W -2b \Delta W_{2} +\frac{\alpha \Delta W}{\alpha  \beta -\gamma ^2}, \\
    0 = & -e W _{(2,2)}+\frac{\beta  W_{2}}{\alpha  \beta -\gamma ^2} - b\Delta^2 W  - 2d\Delta W_{2}+\frac{\gamma \Delta W}{\gamma ^2-\alpha  \beta }, \\
    W  W_{2} = & \frac{\gamma  W_{(2,2)}}{\gamma ^2-\alpha  \beta }-a\Delta^2 W_{2} -b \Delta W_{(2,2)} -b W_{(2,1,2)}-d W_{(2,2,2)}+\frac{\alpha  \Delta W_{2}}{\alpha  \beta -\gamma ^2}, \\
    W = & \frac{\beta  W_{(2,2)}}{\alpha  \beta -\gamma ^2}-b \Delta^2 W_{2}-d\Delta W_{(2,2)}-d W_{(2,1,2)}-e W_{(2,2,2)}+\frac{\gamma \Delta W_{2}}{\gamma ^2-\alpha  \beta },\\
    W  W_{(2,2)} = & \frac{\gamma  W_{(2,2,2)}}{\gamma ^2-\alpha  \beta }-a\Delta^2 W_{(2,2)}-b\Delta W_{(2,2,2)}-b W_{(2,1,2,2)}-d W_{(2,2,2,2)}+\frac{\alpha  \Delta W_{(2,2}}{\alpha  \beta -\gamma ^2}, \\
    W  p_{2}+ W_{2} = & \frac{\beta  W_{(2,2,2)}}{\alpha  \beta -\gamma ^2}-b\Delta^2 W_{(2,2)}-d\Delta W_{(2,2)}-d W_{(2,1,2,2)} -e W_{(2,2,2,2)}+\frac{\gamma \Delta W_{(2,2)}}{\gamma ^2-\alpha  \beta },\\
    W_{2}^2 = & \frac{\alpha  W_{(2,1,2)}}{\alpha  \beta -\gamma ^2}+\frac{\gamma  W_{(2,2,2)}}{\gamma ^2-\alpha  \beta }-a W_{(2,1,1,2)}-2 b W_{(2,1,2,2)}-d W _{(2,2,2,2)},\\
    2 W_{2} = & \frac{\gamma  W_{(2,1,2)}}{\gamma ^2-\alpha  \beta }+\frac{\beta  W_{(2,2,2)}}{\alpha  \beta -\gamma ^2}-b W_{(2,1,1,2)}-2 d W_{(2,1,2,2)}-e W_{(2,2,2,2)}.
\end{align}
%
These are eight independent equations in eight unknowns which may be eliminated to give a single algebraic equation for the resolvent $W(z)$. This spectral curve is degree 4 in both the resolvent and the boundary cosmological constant. We can, at most, tune the parameters of the model to achieve critical behaviour of the form $W(z)\sim(z-z_c)^{4/3}$, which is characteristic of the critical Ising model universality class.


\section{Boundary conditions of the Potts model}

The $q$-state Potts model on random surfaces is defined by the multi-matrix integral 
\beq
\label{pottsModel}
Z_{q}\left(\{t_m\}_{m=2}^{k+2}\right) = \int \prod_{\langle i, j \rangle} e^{N\Tr X_i X_j}\prod_{i=1}^q e^{-N\Tr V_i(X_i)} \rmd X_i
\eeq
where $\langle i,j\rangle$ denotes the product over distinct $i, j$ and we take $V_i(z)  = U_i(z) - z^2/2 \, \forall i$. This is a particular case of the \textit{dilute} $q$-state Potts model which we define to be
\beq
\label{dilutePottsModel}
Z_q\left(t_0, \{t_m\}_{m=2}^{k+2}\right) = \int \rmd Y e^{-N \Tr (t_0 Y^3/3 + Y^2/2)}\prod_{i=1}^q e^{-N\Tr (U_i(X_i)-X_iY)} \rmd X_i.
\eeq
It is straightforward to see that in the case $t_0 = 0$ one can recover \eref{pottsModel} by integrating out $Y$. Taking $U_i(z)=U(z) = \sum_{m=2}^{k+2}t_m z^m/m$, we obtain the $S_q$-symmetric Potts model.

These models are interesting because the cyclical coupling in \eref{pottsModel} prevents one from integrating out the relative angles between the matrices for $q\geqslant 3$. Therefore one is unable to deal with the eigenvalues only and the problem is not amenable to orthogonal polynomial or saddle-point solutions, the usual method of attack for one-matrix models. However, the introduction of the fiducial matrix $Y$, as in \eref{dilutePottsModel}, allows us to disentangle the cyclical coupling and cast the model in the form of an external field problem \cite{gross1, gross2}. The relative angle between the $X_i$ and $Y$ can be integrated out and we may describe the eigenvalues of each $X_i$ in terms of the eigenvalues of $Y$, rendering the problem approachable through saddle-point methods. 

The matrix model studied in the literature thus far typically belong to the class of $S_q$-symmetric Potts models. The Potts matrix model was proposed in \cite{Kazakov1988} wherein the solution was discussed for the case of $q = 1$ and $q = 0$. For $q = 2$ we recover the
Ising model on random surfaces, a two-matrix model, which has been solved \cite{Kazakov1986, Boulatov1987, Eynard:2004mh}. For the particular case of triangulations ($k=1$) and integer $0\leqslant q \leqslant 4$, the analytic properties of the single-matrix resolvent were first computed by Daul \cite{Daul:1994qy} using saddle-point methods. This was extended by Zinn-Justin \cite{ZinnJustin:1999jg}, who first defined the dilute $q$-state Potts matrix model \eref{dilutePottsModel} and was able to derive an algebraic expression for the resolvent function. Using the method of loop equations Bonnet and Eynard \cite{Eynard:1999gp} were able to demonstrate that the resolvent satisfies an algebraic equation for rational $\arccos(q - 2) / \pi $. Atkin, Niedner, and Wheater \cite{Atkin:2015ksy} demonstrated that the resolvents for sums of random matrices also satisfiy algebraic equations for integer $0\leqslant q \leqslant 4$ using saddle-point methods.



%The 3-state Potts model on random triangulated surfaces is defined by a multi-matrix integral with potential
%
%\beq
%V(\mathbf{X}) = \sum_{i=0}^2\frac{t_2 -1}{2} X_i^2 + \frac{t_3}{3} X_i^3 - X_0X_1 -X_1 X_2 - X_2X_0
%\eeq
%
%where we have chosen the coefficients to align with the notation of \cite{Atkin:2015ksy}. This action is manifestly $S_3$-symmetric under exchange of the matrix degrees of freedom.

\subsection{Resolvents as boundary conditions}

It is common to regard resolvents as describing boundary conditions on random graphs. In the diagrammatic expansion of the matrix integral each monomial of order $k+2>2$ generates a $(k+2)$-sided polygon. To leading order in the large $N$ expansion, where planar graphs dominate, we may interpret each term as a tesselation of the sphere by polygons in $q$ equally weighted colours, corresponding to the $q$ different matrices. Given $\sigma\in S_q/(S_p\times S_{q-p})$, we define an associated resolvent function 
%
\begin{eqnarray}
\label{discfct}
	W_{(p|\sigma)}(z)=\frac{1}{N}\left\langle\Tr\frac{1}{z-X_{(p|\sigma)}}\right\rangle\ ,\quad X_{(p|\sigma)}=\sum_{i=1}^pX_{\sigma(i)}\ ,
\end{eqnarray}
%
In the $S_q$-invariant case, for given $p$, all $|S_q/(S_p\times S_{q-p})|={q \choose p}$ partition functions $W_{(p|\sigma)}(z)$ are described by the same function, so that we henceforth abbreviate $W_{(p)}(z):=W_{(p|\sigma_0)}(z)$ for a representative $\sigma_0$. For $p=1$, our definition of $W_{(p)}(z)$ reduces to the one studied in \cite{Kazakov1988,Daul:1994qy,ZinnJustin:1999jg,Bonnet:1999nf,Eynard:1999gp}. These functions generate graphs with a single connected boundary containing $p$ equally weighted colours.

In the statistical lattice model interpretation we regard \eref{pottsModel} as a partition function. Each matrix degree of freedom represents the state of a spin with nearest-neighbour interactions between unequal spins. $S_q$-symmetry under permutations of the matrices renders the $q$ states of the statistical system indistinguishable. We may interpret the resolvent \eref{res} as the disk partition function, generating surfaces with a connected boundary. The matrix, and hence spin, used to define the resolvent defines the set of spins allowed on the boundary. Therefore we term the resolvent of a single matrix the \textit{fixed} boundary condition. Similarly we call the resolvent for the sum of all matrices the \textit{free} boundary condition, and we call the resolvent for the sum of $p$ random matrices the \textit{partially magnetised} or \textit{mixed} boundary condition.

Our goal in the following is to describe resolvents of the 3-state Potts model. These typically obey algebraic relations and so we choose to represent the solution for each resolvent as an algebraic equation $F_{(p)}(x,y)=0$. Asymptotic solutions are presented in appendix A with some high-level overview of the branch structure. We present the algebraic equations up to coefficients containing undetermined low-level trace correlation functions. We substitute these undetermined coefficients with terms $c_{i,j}$ where $i$ and $j$ are the associated exponents of $x$ and $y$, allowing us to describe the essential features of the solution, without being encumbered by the complexity of the full expressions. In practice, one is typically interested in the \textit{critical behaviour} of resolvents as the parameters of the model approach certain values. In 2D dynamical triangulations these correspond to scaling limits of the lattice model where we approach a continuous theory. Detailed knowledge of these coefficients is inessential in determining the critical exponents that describe the universal behaviour of the theory. 

\subsection{3-state Potts model}
\subsubsection{$p=1$}

The potential function is given by
\beq
\label{3Potts_p=1}
V(\mathbf{X}) = \sum_{i=0}^2 \frac{t_2-1}{2}X_i^2 + \frac{t_3}{3}X_i^3 - X_0 X_1 -X_1 X_2 - X_2 X_0.
\eeq
The symmetry group for this model is $S_3\times\mathbb{Z}_2$, corresponding to permutation symmetry of the matrices, and the discrete symmetry $X_i \rightarrow X_i^T \; \forall i$.
With $n_{\mathrm{max}} = 5$ the total number of inequivalent loop equations is 64. We have $\sum_{i=1}^5 |\mathcal{W}_i| = 54$ resolvent functions. This is a closed and consistent system of polynomial equations and so many equations are not independent. Therefore only a subset of these are required to determine the resolvent.



\begin{eqnarray}
F_{(1)}(x,y) & = y^5 + \frac{\left(18 t_2-17\right)}{4 t_3}y^4  -\frac{t_3}{4} x^6 +x^5 \left(\frac{3 t_3 y}{2}+\frac{3 t_2}{2}\right) \nonumber\\
& \quad +x^4 \left(-\frac{1}{4} 13 t_3 y^2+\left(\frac{3}{2}-6 t_2\right) y-\frac{t_3}{2}-\frac{9 t_2^2}{4 t_3}\right) \nonumber\\
& \quad +x^3 \left(3 t_3 y^3+\left(7 t_2-6\right) y^2+\frac{\left(3 t_2^2-6 t_2+2 t_3^2\right) y}{t_3}+2 t_2-\frac{t_2^3}{t_3^2}\right) \nonumber \\
& \quad + x^2 \biggl( -t_3 y^4-\frac{3}{2} \left(t_2-5\right) y^3+\frac{\left(15 t_2^2+54 t_2-10 t_3^2-9\right) y^2}{4 t_3} \nonumber\\ 
& \qquad +\frac{3 \left(5 t_2^3+3 t_2^2-3 t_3^2 t_2+t_3^2\right) y}{2 t_3^2}-\frac{t_3}{4}-\frac{3 t_2^2}{2 t_3}+\frac{3 t_2^4}{t_3^3} \biggr)\nonumber \\
& \quad + x \biggl( \left(-t_2-3\right) y^4+\frac{\left(-9 t_2^2-12 t_2+2 t_3^2+9\right) y^3}{2 t_3} \nonumber \\
& \qquad +\left(-\frac{13 t_2^3}{2 t_3^2}+\left(\frac{9}{2 t_3^2}+2\right) t_2-3\right) y^2 \nonumber \\
& \qquad +\frac{\left(-6 t_2^4+6 t_2^3-6 t_3^2 t_2+t_3^4\right) y}{2 t_3^3}+\frac{t_2}{2}-\frac{t_2^3}{t_3^2} \biggr) \nonumber\\
& \quad + c_{0,3}y^3 + c_{0,2}y^2 + c_{0,1}y +c_{0,0}
\end{eqnarray}


\subsubsection{$p=2$}
In order to compute the planar resolvent for the mixed boundary condition we change the basis of the matrix integral \eref{pottsModel} to $M_0 = X_1 + X_2, M_1 = X_2 + X_0, M_2 = X_0 + X_1$. Under this change of basis the potential becomes 
%
\begin{eqnarray}
\label{mixedAction}
    V(\mathbf{M}) & = &  \sum_{i=0}^2 \frac{3t_2 - 1}{8}  M_i^2  - \frac{t_2 + 1}{4} \left( M_0 M_1 + M_1 M_2 + M_2 M_0 \right) \\
    && \quad + \, \frac{t_3}{8} \big( M_0^2 M_1 + M_0^2 M_2 + M_1^2 M_0 + M_1^2 M_2 \nonumber \\
    && \quad\quad + \, M_2^2 M_0 + M_2^2 M_1 - 3 M_0 M_1 M_2 - 3 M_2 M_1 M_0 \big).  \nonumber
\end{eqnarray}

%
The resulting potential retains the $S_3\times\mathbb{Z}_2$ symmetry of \eref{3Potts_p=1}, but possesses many more interaction terms. The choice of basis is arbitrary, one should be able to reproduce the planar resolvent regardless of how we parameterised the potential. In this case our goal is to compute
%
\begin{equation}
    W_{(2)}(z) = \frac{1}{N} \left\langle \Tr\, \frac{1}{z-M_0} \right\rangle.
\end{equation}
%

In contrast to the fixed-spin boundary condition, we are able to determine a closed set of algebraic constraints on the planar resolvent only if we extend to level $n_{\mathrm{max}} = 6$ loop equations. The increased number of terms in the potential relative to \eref{3Potts_p=1} results in a greater number of resolvent functions present in the loop equations. At this level we have $\sum_{i=1}^6|\mathcal{W}_i|=144$ resolvent functions and $\sum_{i=0}^6|\mathcal{V}_i|=245$ equations. Once redundancies have been removed, these loop equations possess a closed system of equations. We have not determined a minimal set of equations that form closed algebraic system of equations.

The final expression for the spectral curve is degree $10$ in the resolvent $W_{(2)}(z)$ and degree $13$ in the boundary cosmological constant $z$. At first glance, this is rather surprising, as it was shown in \cite{Niedner:2016skw} that the function $G_{(2)}^Y(z)$, which contains the information of the eigenvalue density distribution, satisfies a spectral curve that is degree 5 in $G_{(2)}^Y(z)$ and degree 6 in $z$. However, through detailed analysis of the analytic structure of the two functions one is able to reconcile these observations.

\begin{eqnarray}
F_{(2)}(x,y) &=& y t_3^2 x^{13}+\frac{9}{4} y^2 t_3^2 x^{12}-t_3^2 x^{12}+y \left(13 t_2 t_3-32 t_3\right) x^{12} \nonumber \\
&& +\, 2 y^3 t_3^2 x^{11}+y^2 \left(27 t_2 t_3-88 t_3\right) x^{11} \nonumber \\
&& +\, y \left(\frac{295 t_2^2}{4}-\frac{717 t_2}{2}-2 t_3^2+344\right) x^{11} \nonumber \\
&& +\, \left(\frac{2227 t_2^2}{16}-\frac{7297 t_2}{8}-\frac{3 t_3^2}{2}+\frac{4631}{4}\right) y^2 x^{10} + \frac{7}{8} y^4 t_3^2 x^{10} \nonumber \\
&& +\, y^3 \left(22 t_2 t_3-100 t_3\right) x^{10} + \left(\frac{35 t_2 t_3}{4}-60 t_3\right) y^4 x^9 \nonumber \\
&& +\, \left(\frac{815 t_2^2}{8}-\frac{3805 t_2}{4}-\frac{t_3^2}{2}+\frac{13183}{8}\right) y^3 x^9 \nonumber \\
&& +\, \frac{3}{16} y^5 t_3^2 x^9+\left(\frac{27 t_2 t_3}{16}-20 t_3\right) y^5 x^8 \nonumber \\
&& +\, \left(\frac{2295 t_2^2}{64}-\frac{16605 t_2}{32}-\frac{t_3^2}{16}+\frac{82239}{64}\right) y^4 x^8 \nonumber \\
&& +\, \frac{1}{64} y^6 t_3^2 x^8+\left(\frac{t_2 t_3}{8}-\frac{7 t_3}{2}\right) y^6 x^7 + y^5 \left(6 t_2^2-\frac{311 t_2}{2}+\frac{1185}{2}\right) x^7 \nonumber \\
&& +\, \left(\frac{3 t_2^2}{8}-\frac{193 t_2}{8}+\frac{1291}{8}\right) y^6 x^6 - \frac{1}{4} y^7 t_3 x^6+\left(24-\frac{3 t_2}{2}\right) y^7 x^5 \nonumber \\
&& +\, y^6 \left(\frac{t_2^3}{2 t_3}-\frac{243 t_2^2}{4 t_3}+\frac{3621 t_2}{4 t_3}-\frac{2387}{t_3}\right) x^5 + \frac{3 y^8 x^4}{2}\nonumber \\
&& +\, y^7 \left(-\frac{3 t_2^2}{t_3}+\frac{231 t_2}{2 t_3}-\frac{1109}{2 t_3}\right) x^4 + y^8 \left(\frac{6 t_2}{t_3}-\frac{72}{t_3}\right) x^3 \nonumber \\
&& +\, y^7 \left(-\frac{2 t_2^3}{t_3^2}+\frac{174 t_2^2}{t_3^2}-\frac{1930 t_2}{t_3^2}+\frac{3360}{t_3^2}\right) x^3 \nonumber \\
&& +\, y^8 \left(\frac{6 t_2^2}{t_3^2}-\frac{198 t_2}{t_3^2}+\frac{690}{t_3^2}\right) x^2 \frac{4 y^9 x^2}{t_3} + y^9 \left(\frac{80}{t_3^2}-\frac{8 t_2}{t_3^2}\right) x \nonumber \\
&& +\, y^8 \left(-\frac{108 t_2^2}{t_3^3}+\frac{948 t_2}{t_3^3}+\frac{144}{t_3^3}\right) x + \frac{4 y^{10}}{t_3^2}\nonumber \\
&& +\, \frac{8 y^9 \left(7 t_2+1\right)}{t_3^3} + \sum_{i=1}^{12} \sum_{j=1}^{i} x^{12-i} y^{i-j} c_{i,j}
\end{eqnarray}




\subsubsection{$p=3$}

Let us choose a basis for the potential \eref{3Potts_p=1} given by $M_0 = X_0 + X_1 + X_2, M_1 = X_0 - X_1 + X_2, M_2 = X_0 + X_1- X_2$. Then computing the resolvent corresponding to the free boundary condition, $W_{(3)}(z)$, is equivalent to computing the planar resolvent of the matrix $M_0$ with the following potential,
%
\begin{eqnarray}
\label{freeAction}
V(\mathbf{M}) & = & \frac{(t_2-2)}{4}M_0^2 + \frac{t_2}{4}(M_1^2 + M_2^2) + \frac{t_2}{4}(- M_0 M_1 + M_1 M_2 - M_2 M_0) \nonumber \\
&& + \, \frac{t_3}{8}(M_1^2 M_0 + M_2^2 M_0 + M_1^2 M_2 + M_2^2 M_1)  \\
&& + \, \frac{t_3}{8}(- M_0^2 M_1 - M_0^2 M_2) + \frac{t_3}{12} M_0^3. \nonumber
\end{eqnarray}
%
The symmetry group of \eref{freeAction} possesses a $\mathbb{Z}_2\times\mathbb{Z}_2$ symmetry, with one $\mathbb{Z}_2$ symmetry relating $M_1$ and $M_2$. However, the correlators also possesses a further set of constraints, owing to their origin in the 3-state Potts model, that are not manifest in \eref{freeAction}. These are given by
%
\begin{equation}
    \label{extraCon}
    W_{1}(z) = \frac{1}{3} \Delta W_{(3)} (z),
\end{equation}
%
which is easily deduced by rewriting the terms in the original basis and imposing the full $S_3$ symmetry in the expectation value. 

As in the case of the mixed boundary condition, one must extend to level 6 to find a complete set of constraints. Applying the elimination procedure and \eref{extraCon}, one can determine the spectral curve for the free spin resolvent, given by an algebraic equation that is degree 5 in $z$ and degree 5 in the resolvent $W_{(3)}(z)$. Furthermore, with the relation
%
\begin{equation}
    \label{strange}
    G_{(3)}^Y(z) = z + W_{(3)}(z),
\end{equation}
%
one recovers the spectral curve for $G_{(3)}^Y(z)$ of \cite{Atkin:2015ksy}

\begin{eqnarray}
 F_{(3)}(x,y) &=&  x^6 + 6 x^5 y + 9 x^4 y^2 - 4 x^3 y^3 - 12 x^2 y^4 + \frac{108 y^5}{t_3} \nonumber \\
  && +\,  \frac{27 (26 t_2 - 9) y^4}{t_3^2} + \frac{54 t_2 (6 t_2^2 + t_3^2) x}{t_3^3} \nonumber \\
  && +\, \frac{18 (-18 t_2^4 - 54 t_2^3 + 12 t_3^2 t_2^2 - 18 t_3^2 t_2 + t_3^4) x y}{t_3^4} \nonumber \\
  && -\, \frac{54 (9 t_2^3 + 12 t_2^2 - 9 t_2 + 2 t_3^2) x y^2}{t_3^3} - \frac{6 (39 t_2^2 + 2 t_3^2 - 27) x y^3}{t_3^2}\nonumber \\
  && -\, \frac{36 (t_2 - 1) x y^4}{t_3} + \left( \frac{324 t_2^4}{t_3^4} + \frac{270 t_2^2}{t_3^2} + 9 \right) x^2 \nonumber \\
  && +\, \frac{54 (3 t_2^3 - 15 t_2^2 + 3 t_3^2 t_2 - t_3^2) x^2 y}{t_3^3}  \\
  && +\, \frac{9 (-15 t_2^2 - 54 t_2 + 2 t_3^2 + 9) x^2 y^2}{t_3^2} \nonumber \\
  && -\, \frac{18 (5 t_2 + 3) x^2 y^3}{t_3} + \frac{36 (9 t_2^3 + 2 t_3^2 t_2) x^3}{t_3^3} \nonumber \\
  && +\, \frac{12 (21 t_2^2 - 18 t_2 + 2 t_3^2) x^3 y}{t_3^2} + \frac{36 (t_2 - 2) x^3 y^2}{t_3}\nonumber \\
  && +\, \left( \frac{117 t_2^2}{t_3^2} + 6 \right) x^4 + \frac{18 (4 t_2 - 1) x^4 y}{t_3} \nonumber \\
  && +\, \frac{18 t_2 x^5}{t_3} + y^3 c_{1,4} + y^2 c_{1,3} + y c_{1,2} + c_{1,1}\nonumber 
\end{eqnarray}

\subsection{Dilute 3-state Potts}

\subsubsection{Vacant boundary}

\subsubsection{$p=1$}

\subsubsection{$p=2$}

\subsubsection{$p=3$}


\subsection{General 3-state Potts}
\beq
\label{3Potts_p=1}
V(\mathbf{X}) = \sum_{i=0}^2 \frac{t_{2,i}-1}{2}X_i^2 + \frac{t_{3,i}}{3}X_i^3 - X_0 X_1 -X_1 X_2 - X_2 X_0
\eeq
In this case we have broken the $S_3$ symmetry.

\begin{eqnarray}
F_{(1)}(x,y) & = y^7+y^6 \left(\frac{3 r_2-2}{r_3}+\frac{3 s_2-2}{s_3}+\frac{2 t_2-1}{4 t_3}\right) \nonumber\\
& -\frac{t_3}{4} x^8+\left(2 y t_3+x^7\left(\frac{r_2}{r_3}+\frac{s_2}{s_3}\right) t_3\right) \nonumber \\
& + x^6\biggl(-\frac{13t_3}{2} y^2+\frac{1}{2} \left(t_2+\frac{\left(1-13 r_2\right) t_3}{r_3}+\frac{\left(1-13 s_2\right) t_3}{s_3}+1\right) y \nonumber \\
& \quad +\frac{1}{4} \left(-\frac{6 t_3 r_2^2}{r_3^2}+\frac{\left(s_3 t_2-13 s_2 t_3\right) r_2}{r_3 s_3}+\frac{s_2 t_2}{s_3}-2 t_3-\frac{6 s_2^2 t_3}{s_3^2}+\frac{t_2^2}{t_3}\right)\biggr) \nonumber\\
&  + x^5\biggl(11 t_3 y^3+\frac{3}{2} \left(-2 t_2+\frac{\left(11 r_2-2\right) t_3}{r_3}+\frac{\left(11 s_2-2\right) t_3}{s_3}-2\right) y^2 \nonumber \\
& \quad +\frac{3}{2} \biggl(-\frac{t_2^2}{t_3}+\frac{r_2 \left(5 r_2-1\right) t_3}{r_3^2}+\frac{s_2 \left(5 s_2-1\right) t_3}{s_3^2}+2 t_3 \nonumber \\
& \qquad -\frac{s_2 \left(2 t_2+1\right)}{s_3}-\frac{s_2 t_3+r_2 \left(s_3 \left(2 t_2+1\right)+\left(1-11 s_2\right) t_3\right)}{r_3 s_3}\biggr) y \nonumber \\
& \qquad -\frac{15 r_2^2 s_2 t_3}{4 r_3^2 s_3}+\frac{15 r_2 s_2^2 t_3}{4 r_3 s_3^2}-\frac{3 r_2 s_2 t_2}{2 r_3 s_3}+\frac{r_2^3 t_3}{r_3^3}-\frac{3 r_2^2 t_2}{4 r_3^2}+\frac{3 r_2 t_3}{2 r_3} \nonumber \\
& \qquad -\frac{3 r_2 t_2^2}{4 r_3 t_3}+\frac{3 s_2 t_3}{2 s_3}+\frac{s_2^3 t_3}{s_3^3}-\frac{3 s_2^2 t_2}{4 s_3^2}-\frac{3 s_2 t_2^2}{4 s_3 t_3} \biggr) \nonumber\\
&  + x^4 \biggl(-\frac{41t_3}{4}y^4+\frac{1}{2} \left(14 t_2+\frac{\left(13-41 r_2\right) t_3}{r_3}+\frac{\left(13-41 s_2\right) t_3}{s_3}+16\right) y^3 \nonumber \\
& \quad +\frac{1}{4} \biggl(\frac{s_2 \left(42 t_2+32\right)-4 \left(t_2+1\right)}{s_3}+\frac{\left(-55 r_2^2+26 r_2-1\right) t_3}{r_3^2} \nonumber \\
& \qquad +\frac{\left(-55 s_2^2+26 s_2-1\right) t_3}{s_3^2}+\frac{s_3 \left(r_2 \left(42 t_2+32\right)-4 \left(t_2+1\right)\right)+\left(26 s_2+r_2 \left(26-123 s_2\right)+2\right) t_3}{r_3 s_3}+\frac{14 t_2^2+2 t_2-28 t_3^2-1}{t_3}\biggr) y^2+c_{5,2} y+c_{5,1}\biggr) 
& \quad +c_{1,1}+y c_{1,2}+y^2 c_{1,3}+y^3 c_{1,4}+y^4 c_{1,5}+y^5 c_{1,6}
\end{eqnarray}
